\section{Quantum Bayesian Networks (\Qfactor-based)}
\label{sec:qbn}

We define quantum Bayesian networks in terms of \qfactor s.
Our definition is equivalent to the instrument-based definition in~\cite{GBN2014}, however moving to \qfactor s allows the semantics to address compositionality of sub-networks.
To each node, we associate a special kind of \qfactor, called a \qCPT.
This ensures that the semantics of the full network is exactly a probability distribution, and not any arbitrary function with values in $\nnReal$.

\subparagraph{Quantum CPTs.}

We define a class of \Qfactor s which plays a similar role to that of CPTs.

\begin{definition}[\QCPT\ (\qCPT)]
\label{def:condition_qfactor_qbn}
A \qfactor\ $\someqfact$ over $(\bX,\qreg)$ is
\begin{itemize}
\item
a \definitive{\qCPT\ for $Z \in \bX$ given $(\bY,\qreg)$} if $\bX = \{Z\} \uplus \bY$, and $\forall \by\in\setofrv{\bY}$:
\begin{equation*}
\left(\sum_{Z} \someqfact\right)(\by)
=
\id[\L(\HH_\qreg)]
\end{equation*}
\item
a \definitive{\qCPT\ for $\qreg' \subseteq \qreg$ given $(\bX, \qreg'')$} if $\qreg = \qreg' \uplus \qreg''$, and $\forall \bx\in\setofrv{\bX}$:
\begin{equation*}
\left(\sum_{\qreg'} \someqfact\right)(\bx)
=
\id[\L(\HH_{\qreg''})]
\end{equation*}
\end{itemize}
\end{definition}

The two limit cases correspond to standard, familiar notions.
\begin{itemize}
\item
In the \emph{pure classical case} ($\qreg=\emptyset$), being a \qCPT\ amounts to being a CPT.
Indeed, a \qCPT\ for $X$ given $(\bY, \emptyset)$ is a \qfactor\ $\someqfact$ over $(\{X\}\uplus\bY,\emptyset)$ which satisfies $\sum_{x\in\setofrv{X}} \someqfact(x,\by) = 1$.
\item
In the \emph{pure quantum case} ($\bX=\emptyset$), being a \qCPT\ amounts to being \emph{Choi state}, a well-known condition on operators that we will detail in \cref{sec:choi_iso}.
\end{itemize}

\subparagraph{Quantum Bayesian Networks.}

We define quantum Bayesian networks as DAGs with a \qCPT for each node.

\begin{definition}[Quantum Bayesian Network (\qqbn)]
A \definitive{quantum Bayesian network} over a set $\bX$ of \rvs\ and a set $\bQ$ of disjoint non-empty quantum registers is a pair $(\G,\someqfactset)$ where:
\begin{itemize}
\item
$\G$ is a directed acyclic graph over the set of nodes $\bX\cup\bQ$;
\item
each out-edge of a node $\qreg\in\bQ$ is labelled by a non-empty register $\qreg_i$, with $\uplus_i \qreg_i = \qreg$;
\item
each out-edge of a node $X\in\bX$ is labelled by $X$;
\item
$\someqfactset$ assigns to each $V \in \bX\cup\bQ$ a \qCPT\ for $V$ given the \rvs\ and qubits labelling its in-edges.
\end{itemize}
\end{definition}

\begin{example}[Alice's \qCPT]
Consider \cref{fig:Bell} with the setting of \cref{ex:Alice_instrument}.
The node $A$ is given the \qfactor\ $\someqfact^A: \setofrv{A}\times\setofrv{X} \to \L(\HH_{q_1})$ below, that is a \qCPT\ for $A$ given $(\{X\},\{q_1\})$:

\begin{adjustbox}{}
\begin{tabular}{ccc@{ }c@{ }l@{}l@{\qquad\qquad}ccc@{ }c@{ }l@{}l}
$(a^t,x^t)$ & $\mapsto$ & $\ket{0}\bra{0}$ & $=$ &
&
$\left[\begin{smallmatrix}
1 & 0\\
0 & 0\\
\end{smallmatrix}\right]$
&
$(a^f,x^t)$ & $\mapsto$ & $\ket{1}\bra{1}$ & $=$ &
&
$\left[\begin{smallmatrix}
0 & 0\\
0 & 1\\
\end{smallmatrix}\right]$
\\
$(a^t,x^f)$ & $\mapsto$ & $\ket{+}\bra{+}$ & $=$ &
$\frac{1}{2}$
&
$\left[\begin{smallmatrix}
1 & 1\\
1 & 1\\
\end{smallmatrix}\right]$
&
$(a^f,x^f)$ & $\mapsto$ & $\ket{-}\bra{-}$ & $=$ &
$\frac{1}{2}$
&
$\left[\begin{smallmatrix}
1 & -1\\
-1 & 1\\
\end{smallmatrix}\right]$
\end{tabular}
\end{adjustbox}
\end{example}

\subsection{Semantics and Compositionality}

The semantics of a \qqbn\ is a probability distribution, defined akin to that of a BN (\cref{lem:BNsem}).

\begin{definition}
For $\mathcal{B} = (\G,\someqfactset)$ a \qqbn\ over \rvs\ $\bX$ and registers $\bQ$, its \definitive{semantics} is:
\begin{equation*}
\sem{\mathcal{B}} \defeq \bigodot_{V \in \bX\cup\bQ} \someqfactset(V)
\end{equation*}
\end{definition}

\begin{proposition}
\label{lem:sem_qbn_dist}
With the above notations, $\sem{\mathcal{B}}$ is a probability distribution over $\bX$.
\end{proposition}
The proof is obtained in a similar way to that for the classical result~\cite{Pearl86}.

The formalism allows us to naturally interpret sub-components of a \qqbn, as in \cref{fig:compositional}.
Call $(\G',\someqfactset')$ a \definitive{sub-network} of a \qqbn\ $(\G,\someqfactset)$ when $\G'$ is a sub-graph of $\G$ and $\someqfactset'$ is the restriction of $\someqfactset$ to nodes of $\G'$.
%\footnote{This corresponds to the notion of conditional Bayesian networks~\cite[Section~5.6]{KollerBook}.}
The semantics of a sub-network is defined as for the full network: the product of its \qfactor s.
Compositionality follows from associativity of $\compfact$.

\begin{proposition}[Compositionality]
\label{lem:compositionality_qbn}
Let $\mathcal{B} = (\G,\someqfactset)$ be a \qqbn, and $\mathcal{B}_1 = (\G_1,\someqfactset_1)$ and $\mathcal{B}_2 = (\G_2,\someqfactset_2)$ a partition in two sub-networks (\ie\ $\G = \G_1 \uplus \G_2$).
Then:
\begin{equation*}
\sem{\mathcal{B}} = \sem{\mathcal{B}_1} \compfact \sem{\mathcal{B}_2}
\end{equation*}
\end{proposition}

% \begin{remark}
% Observe that, in a sub-network $(\G',\someqfactset')$, some \rv\ or some register can appear in the scope of a \qfactor\ of $\someqfactset'$ while not belonging to the nodes of $\G'$, as opposed to what happens in a quantum Bayesian network.
% \end{remark}

\subsection{Equivalence with Instrument-based Quantum Bayesian Networks}
\label{sec:choi_iso}

\subparagraph{Choi states and \qCPTs.}

The \definitive{Choi-Jamiolkowski isomorphism} (see \eg~\cite{Watrous_book}) is a classical, well-known result allowing to see linear maps from $\L(\HH_1)$ to $\L(\HH_2)$ as matrices in $\L(\HH_1 \otimes \HH_2)$.
It has many applications, basic ones being to check easily whether two quantum operations are equal---by checking their images are the same---or whether a map is completely positive; we refer to~\cite{Watrous_book} for details.
Let us denote by $\choiiso$ the Choi-Jamiolkowski isomorphism.
Given a quantum operation $\Eop : \L(\HH_1) \to \L(\HH_2)$, complete positivity of $\Eop$ is equivalent to positivity of $\choiiso(\Eop)$, and trace-preserving of $\Eop$ is equivalent to $\ptr{\HH_2}(\choiiso(\Eop)) = \id[\L(\HH_1)]$.
Hence, operators $\M\in\L(\HH_1 \otimes \HH_2)$ that are positive and such that $\ptr{\HH_2}(\M) = \id[\L(\HH_1)]$ are exactly the images by $\choiiso$ of quantum operations, and are called \definitive{Choi states}.
Please notice that being a \qCPT\ for a quantum register $Q'$, \ie\ the second item of \cref{def:condition_qfactor_qbn}, corresponds to being a Choi state (once values for the \rvs\ are fixed) with $\left(\sum_{\qreg'} \someqfact\right)(\bx) = \ptr{\HH_{\qreg'}}(\someqfact(\bx)) = \id[\L(\HH_{\qreg''})]$---the positivity requirement being already part of the definition of a \Qfactor.

\subparagraph{Equivalence with Instrument-based Quantum Bayesian Networks.}

Our definition of %\qfactor-based 
quantum Bayesian networks is equivalent to the instrument-based one of~\cite{GBN2014}.
Indeed, the definition of the DAGs is the same, and it is straightforward to provide a bijection between a \qCPT associated to a node in our formulation, and instruments associated to the same node in~\cite{GBN2014}.
The proof is similar to that for the Choi-Jamiolkowski isomorphism (see \cref{app:equivalence_qbns}).
The composition of instruments in~\cite{GBN2014} corresponds to the product of \qCPTs.
